\section{16. T1--T8: aktueller
Gesamtstatus}\label{t1t8-aktueller-gesamtstatus}

{\def\LTcaptype{none} % do not increment counter
\begin{longtable}[]{@{}
  >{\raggedright\arraybackslash}p{(\linewidth - 4\tabcolsep) * \real{0.3333}}
  >{\raggedright\arraybackslash}p{(\linewidth - 4\tabcolsep) * \real{0.3333}}
  >{\raggedright\arraybackslash}p{(\linewidth - 4\tabcolsep) * \real{0.3333}}@{}}
\toprule\noalign{}
\begin{minipage}[b]{\linewidth}\raggedright
Tor
\end{minipage} & \begin{minipage}[b]{\linewidth}\raggedright
Vorliegende Bausteine
\end{minipage} & \begin{minipage}[b]{\linewidth}\raggedright
Fehlender vollständiger Nachweis
\end{minipage} \\
\midrule\noalign{}
\endhead
\bottomrule\noalign{}
\endlastfoot
T1 -- Ursprung und Auswahl & Markierter Compiler, konkrete Tensoren,
Clock in Spin(10), Unterbestimmtheitszeuge & Auswahl des tatsächlichen
Prozesskerns, der Primitive und der physisch unterscheidbaren
Parameter \\
T2 -- Half-Charge und E8-Feld & Endliche Typen, geladene native Antwort,
abstrakte Nahtkandidaten & Derselbe renormierte Half-Charge-Adapter mit
Energie-, Adjungierten-, E8- und Clockkontrolle \\
T3 -- gemeinsamer 3+1D-Parent & Bedingte Transfermodelle, neue
Ein-Cartan-Referenz, operative Wirkungskriterien & Ursprüngliche
räumliche Komposition, lokale Zugriffe und gemeinsame 3+1D-Dynamik \\
T4 -- chirale Materie & Korrigierte Feldkanäle und vollständige innere
Bilineare & Positiver chiraler Feldadapter, Maß, Anomalien und
kontrollierte Spiegelentkopplung \\
T5 -- Kontinuum und Dynamik & Exakte endliche Sektoren, H-Kette, Lücken
und Adapterfehlervertrag & Größenuniforme Kontrolle, relativistischer
wechselwirkender Grenzwert, Clustering und Streuung \\
T6 -- Parameter und Spektren & Compilerformeln, datierte Vergleiche,
explizit verschiedene \(g/\Delta\)-Antworten & Gemeinsamer physischer
Transfer zu Kopplungen, Familienmassen, Neutrinos und Skala \\
T7 -- Gravitation & Lorentz-Kegelansicht, bedingte geometrische
Strukturen & Dynamischer positiver masseloser Spin 2 und universelle
konsistente Kopplung aus demselben Parent \\
T8 -- Zustand und Instrumente & Modellgrundzustand, bedingte Rampe,
Record-Gram, Referenzbilanz & Ursprüngliche Zustandsregel, physische
Präparation, Apparate, Records und deren Ressourcen \\
\end{longtable}
}

\textbf{Alle acht vollständigen Tore bleiben offen.} Das schränkt die
exakten Teilresultate nicht ein; es verhindert ihre unzulässige
Verallgemeinerung. Die vorliegende Arbeit beansprucht keine vollständige
TOE, keine experimentelle Bestätigung und keinen globalen mathematischen
Durchbruch.

\section{17. Schluss}\label{schluss}

Der neueste TFPT-/Universalraum-Stand ist mathematisch substanzieller
und strukturell klarer als eine Sammlung von Analogien. Der markierte
Matrixkern, die native Paarregel, die Grundzustandsantwort und die
korrigierte Singulett-/Lanczos-Struktur sind konkrete Gegenstände. Die
Fortschritte verbinden jedoch bislang verschiedene Ebenen unter
zusätzlichen Verträgen.

Der eigene Lösungsversuch zeigt drei tragfähige Verbindungen: Ein
positiver Wortkern liefert einen gemeinsamen Prozessraum und ein
strenges Äquivalenzkriterium. Eine explizite kleine Referenz macht einen
konkreten Transfer mit erhaltener Gesamtladung möglich und bezahlt dabei
den Referenzwechsel. Eine physische Projektorfamilie kann Zustandsmetrik
und geometrische Phase aus demselben Tangentialprodukt erzeugen.
Gemeinsam präzisieren sie, wie Herkunft, Aufzeichnung, Transport und
Geometrie in einer einzigen Ausführung behandelt werden könnten.

Die alles verbindende physische Lösung ist damit nicht gefunden. Der
zentrale offene Inhalt ist aber jetzt deutlich: Die ursprüngliche
TFPT-Regel muss einen solchen gemeinsamen Prozess samt Zustand und
lokalen Zugriffen auswählen. Erst wenn diese Auswahl gelingt und
derselbe Prozess einen kontrollierten relativistischen Grenzwert
besitzt, können die bereits vorhandenen algebraischen und spektralen
Ergebnisse zu einer vollständigen Theorie unserer Welt zusammenwachsen.

\section{Anhang A. Notation und
Modellgrenzen}\label{anhang-a.-notation-und-modellgrenzen}

{\def\LTcaptype{none} % do not increment counter
\begin{longtable}[]{@{}
  >{\raggedright\arraybackslash}p{(\linewidth - 2\tabcolsep) * \real{0.5000}}
  >{\raggedright\arraybackslash}p{(\linewidth - 2\tabcolsep) * \real{0.5000}}@{}}
\toprule\noalign{}
\begin{minipage}[b]{\linewidth}\raggedright
Symbol
\end{minipage} & \begin{minipage}[b]{\linewidth}\raggedright
Bedeutung
\end{minipage} \\
\midrule\noalign{}
\endhead
\bottomrule\noalign{}
\endlastfoot
\(\Omega_4\) & antisymmetrisches Vierer-Singulett im Tetramermodell \\
\(\Omega\) & übernommener nativer Grundzustand bei \(N=64\) \\
\(F\) & voll besetzte native Fermionreferenz ohne Bosonen \\
\(N\) & native Teilchenladung \(N_f+2N_b\) \\
\(Q_S,Q_R\) & einzelne innere Cartanladung des Transferzeugen und
Referenzladung; nicht \(N\) \\
\(g,\Delta\) & native Paarstärke und Bosonenergie \\
\(J\) & je nach ausdrücklich benanntem Modell Tetramerstärke oder
zusätzlicher Mischer; keine automatische Identifikation \\
\(S=T^\dagger T\) & Gramoperator eines Adapters \\
\(\mathsf S\) & Swapoperator; verschieden vom Gramoperator \(S\) \\
\(P\) & orthogonaler Projektor; im Matrixordnungsabschnitt steht \(P\)
stattdessen für die dort angegebene Basismatrix \\
\(\Gamma(u,v)\) & vollständiger Wort-Gramkern; nicht die Exteriorhebung
\(\Gamma(p)\) \\
\(\mu_n\) & Energiemomente der Referenz \(F\) \\
\(m_n\) & normierte Momente der geladenen nativen Antwort auf
\(\Omega\) \\
\end{longtable}
}

Die Notationsüberschneidungen stammen teils aus den Quellen. In jedem
Abschnitt wird der jeweilige Vertrag ausdrücklich angegeben. Zahlen
werden nicht zwischen diesen Verträgen übertragen.

\section{Anhang B. Reproduzierbarkeit und Grenzen der neuen
Prüfung}\label{anhang-b.-reproduzierbarkeit-und-grenzen-der-neuen-pruxfcfung}

Das begleitende Prüfpaket enthält den neuen Prüfer, den tatsächlich
verwendeten Tensor, beide Ergebnisdateien, den isolierten
H-Ketten-Replay, die Diagrammquellen, das Dokumentmanifest und die
eingefrorenen Quellen. Die 57 Bedingungen zählen Prüfbedingungen, keine
57 unabhängigen Theoreme. Die analytischen Beweise in den Abschnitten
8--10 sind Bestandteil der wissenschaftlichen Begründung; sie sind nicht
in einem formalen Beweisassistenten verifiziert.

Die Tests prüfen insbesondere: die vollständige achtdimensionale
Referenzmatrix und ihre invariante Einschränkung; die
Ladungskommutatoren; den endlichen geometrischen Phasenzeugen samt
Wahrscheinlichkeit; die glatte Projektorkrümmung und Metrik; Record-Gram
und Gegenkontrolle; die symbolischen Momentenbeziehungen und rationale
Ritz-Einschließung; den Hash, die ganzzahligen Einträge und
Gramidentität von W sowie das vollständige charakteristische Polynom
seines Unterstützungsgraphen.

Während der Entwicklung wies ein erster Vergleich eine korrekt
äquivalente trigonometrische Form nicht als identisch aus. Ein anderer
Test verwarf eine zunächst falsch angesetzte Spektralmultiplizität. Die
endgültige Rechnung verwendet die exakt faktorisierte Charakteristik des
tatsächlichen Tensors und erhält die Laplacelücke zwölf. Beide Fehler
wurden vor der Auslieferung korrigiert; fehlgeschlagene
Entwicklungsläufe zählen nicht als erfolgreiche Verifikation.

Die neue Referenzrechnung prüft die blockweise Identität mit H4. Sie
übernimmt den rigorosen Transfer-Endpunktbeweis aus S033. Die Abbildung
zeigt eine ergänzende numerische Zeitkurve. Die nativen Grundzustands-,
Polgewichts- und Restkoeffizientensätze stammen aus den bezeichneten
früheren Nachweisen. Die jetzige Rechnung beweist weder die vollständige
historische Enumeration erneut noch die Herkunft ihrer physikalischen
Annahmen.

Die PDFs wurden aus diesem Markdown über LaTeX erzeugt; Formeln bleiben
Text und Vektorgrafik. Das Auslieferungsprotokoll dokumentiert
Seitenzahl, Textabdeckung, Dateihashes und visuelle Stichproben. Das
vollständige Quellenmanifest unterscheidet Eingaben und neue Ergebnisse.
Der Quellenstand ist eine eingefrorene Forschungsaufnahme, kein
Versprechen, nach Abschluss eintreffende parallele Änderungen bereits zu
enthalten.

\section{Anhang C. Zentrale lokale
Quellen}\label{anhang-c.-zentrale-lokale-quellen}

Die vollständigen Pfade, Hashes und Zeitstempel aller inventarisierten
Dokumente stehen im Quellenmanifest. Die Kennungen im Haupttext beziehen
sich auf die folgenden eingefrorenen Dateien:

\begin{itemize}
\tightlist
\item
  \textbf{S001:} \emph{TFPT / Universalraum: Präparation, Restantwort
  und die Herkunft räumlicher Teile. Fundamentale Fortsetzung v1.6.6},
  15.09.2026. §§2--5: Präparation, Restantwort, Schnittschranke,
  Kinetik.
\item
  \textbf{S003:} \emph{TFPT / Universalraum: die fundamentale
  Reduktion}, 15.09.2026. §§3--8: volle symmetrische Antwort,
  Multiplizität, Zustandswahl und Interventionskriterium.
\item
  \textbf{S033:} \emph{Clock, gemeinsame Quelle und tatsächlicher
  Transport. Konsolidierung v1.6.6}, 15.09.2026. §§1--7: native
  Intervalle, Clock, Casimire, H4, Feldadapter und Überlappung.
\item
  \textbf{S046:} \emph{TFPT / Universalraum: Erklärung und
  Rekonstruktion. Hauptdokument v1.6}, 15.09.2026, 168 PDF-Seiten. Kap.
  3--9: Grundstruktur; Kap. 27--33: physikalische Anschlüsse; Kap.
  34--37: Matrixkern, native Dynamik und Ursprung.
\item
  \textbf{S068:} \emph{Die richtige einfache Kette. Vollständige
  Konsolidierung v1.6.7}, Ergebnisbericht, 15.09.2026. §§3--8: vier
  Singuletts, richtige H-Kette, Bilineare und Unterbestimmtheit.
\item
  \textbf{S070:} vollständige Hauptfassung von v1.6.7 einschließlich
  historischer v1.6.6-Anhänge.
\item
  \textbf{S071:} Repositorydokument \emph{What is genuinely open}.
  Zentraler T1--T8-Status; neuere Forschungen werden durch die einzeln
  gepinnten Ergebnisberichte ergänzt.
\item
  \textbf{S073:} Nachgereichter Nutzertext \emph{Operationssatz, nativer
  Grundzustand und Feldwörterbuch}, v1.6.3; byteidentisch mit dem
  bereits gelesenen Repositorybericht. Abgleich und offene
  Forschungsrichtung in §6.4.
\item
  \textbf{S074:} späteste eingefrorene Vollfassung v1.6.7 mit Teil B,
  \emph{Neue Vogelperspektive: derselbe Prozess in verschiedenen
  Ansichten}. B3: Wirkungstest; B4: Z4 und kanonische Reduktion; B5:
  Kodierung und Schatten. Die dort berichteten 822 Prüfbedingungen
  wurden hier nicht vollständig erneut ausgeführt.
\item
  \textbf{S072:} \emph{TFPT Compiler / Universalraum Paper v1.2},
  14.09.2026. Grundannahmen, History-Reparatur, lokale U-Regel und
  physikalische Transfergrenzen.
\end{itemize}

\section{Anhang D. Externe Literatur und
Referenzdaten}\label{anhang-d.-externe-literatur-und-referenzdaten}

Die folgenden Arbeiten liefern etablierte Methoden oder datierte
Vergleichsdaten. Sie beweisen nicht die TFPT-spezifische Herkunft der
jeweiligen Eingaben. URLs wurden während dieser Synthese geprüft; bei
einzelnen Arbeiten wurde nur die offizielle bibliographische Seite
beziehungsweise Zusammenfassung konsultiert, soweit der Haupttext keine
weitergehende Behauptung daraus übernimmt.

\textbf{E1.} G. Chiribella, G. M. D'Ariano, P. Perinotti:
\emph{Theoretical framework for quantum networks}. Phys. Rev.~A 80,
022339 (2009). \href{https://arxiv.org/abs/0904.4483}{Originalarbeit}.

\textbf{E2.} F. A. Pollock et al.: \emph{Operational Markov condition
for quantum processes}. Phys. Rev.~Lett. 120, 040405 (2018).
\href{https://arxiv.org/abs/1801.09811}{Originalarbeit}.

\textbf{E3.} S. Jansen, M.-B. Ruskai, R. Seiler: \emph{Bounds for the
adiabatic approximation with applications to quantum computation}. J.
Math. Phys. 48, 102111 (2007).
\href{https://arxiv.org/abs/quant-ph/0603175}{Originalarbeit}.

\textbf{E4.} S. D. Bartlett, T. Rudolph, R. W. Spekkens: \emph{Reference
frames, superselection rules, and quantum information}. Rev.~Mod. Phys.
79, 555 (2007).
\href{https://arxiv.org/abs/quant-ph/0610030}{Originalarbeit}.

\textbf{E5.} F. Wilczek, A. Zee: \emph{Appearance of Gauge Structure in
Simple Dynamical Systems}. Phys. Rev.~Lett. 52, 2111 (1984).
\href{https://doi.org/10.1103/PhysRevLett.52.2111}{Verlagsquelle}.

\textbf{E6.} S. Bravyi, M. B. Hastings, F. Verstraete:
\emph{Lieb-Robinson bounds and the generation of correlations and
topological quantum order}. Phys. Rev.~Lett. 97, 050401 (2006).
\href{https://arxiv.org/abs/quant-ph/0603121}{Originalarbeit}.

\textbf{E7.} A. Bisio, G. M. D'Ariano, P. Perinotti, A. Tosini:
\emph{Free quantum field theory from quantum cellular automata:
derivation of Weyl, Dirac and Maxwell quantum cellular automata}.
\href{https://arxiv.org/abs/1601.04832}{Originalarbeit}.

\textbf{E8.} P. Nath, R. M. Syed: \emph{Complete Cubic and Quartic
Couplings of 16 and bar 16 in SO(10) Unification}.
\href{https://arxiv.org/abs/hep-th/0109116}{Originalarbeit}.

\textbf{E9.} NIST: \emph{Fundamental Physical Constants --- Complete
Listing, 2022 CODATA adjustment}. Am 15.09.2026 abgerufene datierte
Referenz.
\href{https://physics.nist.gov/cuu/Constants/Table/allascii.txt}{Offizielle
Tabelle}.

\textbf{E10.} S. Bravyi, D. P. DiVincenzo, D. Loss:
\emph{Schrieffer-Wolff transformation for quantum many-body systems}.
Ann. Phys. 326 (2011), 2793--2826. Methodischer Hintergrund für
kontrollierte effektive Dynamik.
\href{https://arxiv.org/abs/1105.0675}{Originalarbeit}.

\textbf{E11.} F. Pastawski, B. Yoshida, D. Harlow, J. Preskill:
\emph{Holographic quantum error-correcting codes: Toy models for the
bulk/boundary correspondence}. JHEP 06 (2015), 149.
\href{https://arxiv.org/abs/1503.06237}{Originalarbeit}.

\textbf{E12.} X. Dong, D. Harlow, A. C. Wall: \emph{Reconstruction of
Bulk Operators within the Entanglement Wedge in Gauge-Gravity Duality}.
Phys. Rev.~Lett. 117, 021601 (2016).
\href{https://arxiv.org/abs/1601.05416}{Originalarbeit}.

\end{document}
